%\documentclass[11pt,a4paper,twoside]{article}

\usepackage[utf8]{inputenc}

\usepackage[T1,T2A]{fontenc}

\usepackage[english.ancient,russian.ancient]{babel}

\usepackage{graphicx}

\usepackage{array}

\usepackage[doi]{samgtu-bib}

\usepackage{samgtu}

\usepackage{samgtu-name}

\usepackage{mathtools}

\mathtoolsset{showonlyrefs}

\usepackage{desclist}

\usepackage{euscript}

\usepackage{mathrsfs} 

\usepackage{multicol}

\usepackage{mathrsfs}

\usepackage{cite}

\usepackage{cmap}

\usepackage{qrcode}

\allowdisplaybreaks[4]

\usepackage[nodayofweek]{datetime}

\usepackage[thinlines]{easybmat}



\renewcommand{\JournalYear}{2018}

\renewcommand{\JournalVolume}{22}

\renewcommand{\JournalNumber}{4}



\newdate{JournalDateSign}{29}{12}{2018}% Дата подписания Вестника в печать

\renewcommand{\JournalDateSign}{\displaydate{JournalDateSign}}

\newdate{JournalDatePrint}{16}{01}{2019}% Дата выхода Вестника из печати

\renewcommand{\JournalDatePrint}{\displaydate{JournalDatePrint}}

%\renewcommand{\JournalUPL}{16.56} % Усл. печ. л.

%\renewcommand{\JournalUIL}{16.53} % Уч.-изд. л.

\renewcommand{\JournalUPL}{15.24} % Усл. печ. л.

\renewcommand{\JournalUIL}{15.21} % Уч.-изд. л.

\renewcommand{\JournalRegNumber}{220/18} % Регистрационный номер

\renewcommand{\JournalOrderNumber}{2} % Номер заказа



%\renewcommand{\JournalRMonth}{Март} % Месяц выхода Вестника

%\renewcommand{\JournalEMonth}{March} % Месяц выхода Вестника

%\renewcommand{\JournalRMonth}{Июнь} % Месяц выхода Вестника

%\renewcommand{\JournalEMonth}{June} % Месяц выхода Вестника

%\renewcommand{\JournalRMonth}{Сентябрь} % Месяц выхода Вестника

%\renewcommand{\JournalEMonth}{September} % Месяц выхода Вестника

\renewcommand{\JournalRMonth}{Декабрь} % Месяц выхода Вестника

\renewcommand{\JournalEMonth}{December} % Месяц выхода Вестника



\newcommand{\totop}[1]{\raisebox{6pt}[1pt][2pt]{#1}}

\newcommand{\tottop}[1]{\raisebox{12pt}[1pt][2pt]{#1}} 

\newcommand{\totttop}[1]{\raisebox{18pt}[1pt][2pt]{#1}} 

\newcommand{\tottttop}[1]{\raisebox{24pt}[1pt][2pt]{#1}} 

\newcommand{\totttttop}[1]{\raisebox{30pt}[1pt][2pt]{#1}} 

\newcommand{\tobot}[1]{\raisebox{-6pt}[1pt][2pt]{#1}} 

\newcommand{\tobbot}[1]{\raisebox{-12pt}[1pt][2pt]{#1}} 

\newcommand{\tobbbot}[1]{\raisebox{-18pt}[1pt][2pt]{#1}}



\graphicspath{{./00PIC/}}

%

\vsgtu{1581} %registration number

\FirstPage{684}

\LastPage{698}



%

\ResearchArticle

%

%\pdfcrop

%%\draft



%\newcommand{\ruipm}{\ruaffid{755}{Институт проблем механики им.~А.\,Ю.~Ишлинского РАН,\\

%Россия, 119526, Москва, просп. Вернадского, 101, корп.~1}}

%\newcommand{\enipm}{\enaffid{755}{A.~Ishlinsky Institite for Problems in Mechanics, Russian Academy of Sciences, \\ 101, pr. Vernadskogo, Moscow, 119526, Russian Federation}}

%

%\newcommand{\rursuh}{\ruaffid{1416}{Российский государственный гуманитарный университет,\\ 

%Россия, 125993, Москва, Миусская площадь, 6}}

%\newcommand{\enrsuh}{\enaffid{1416}{Russian State University for the Humanities,\\ 

%9,  Miusskaya pl., Moscow, 125993, Russian Federation}}

%

%\newcommand{\rumephi}{\ruaffid{7201}{Национальный исследовательский ядерный университет «МИФИ»,\\ 

%Россия, 115409, Москва, Каширское ш., 31}}

%\newcommand{\enmephi}{\enaffid{7201}{National Engineering Physics Institute ``MEPhI'',\\ 

%31,  Kashirskoe shosse, Moscow, 115409 Russian Federation}}

%

%\newcommand{\rubauman}{\ruaffid{1019}{Московский государственный технический университет им.~Н.~Э.~Баумана, \\ Россия, 105005, Москва, 2-я Бауманская ул., 5/1}}

%\newcommand{\enbauman}{\enaffid{1019}{N.~E.~Bauman Moscow State Technical University, \\ 5/1,  2-ya Baumanskaya st., Moscow, 105005, Russian Federation}}



%\usepackage{pgf,pgfplots,tikz}

%\usetikzlibrary{arrows,snakes,backgrounds,patterns,shapes.geometric,calc}

%

%\begin{document}



%

%\hfill \qrcode[hyperlink,height=.8in]{http://doi.org/10.14498/vsgtu1581}

%\vspace{-.8in}



\udc{517.958:539.3(1); 539.385}

\msc{74D99}



\rutitle{Кручение растущего вала}

\rutitleshort{Кручение растущего вала}



\entitle{Torsion of a growing shaft}

\entitleshort{Torsion of a growing shaft}



\ruauthor{А.~В.~Манжиров\affil{1,2,3}, М.~Н.~Михин\affil{4}, \\ Е.~В.~Мурашкин\affil{1,2,3}}{М\,а\,н\,ж\,и\,р\,о\,в~А.~В., М\,и\,х\,и\,н~М.~Н., М\,у\,р\,а\,ш\,к\,и\,н~Е.~В.}

\enauthor{A.~V.~Manzhirov\affil{1,2,3}, M.~N.~Mikhin\affil{4},  E.~V.~Murashkin\affil{1,2,3}}{M\,a\,n\,z\,h\,i\,r\,o\,v~A.~V., M\,i\,k\,h\,i\,n~M.~N., M\,u\,r\,a\,s\,h\,k\,i\,n~E.~V.}



\ruaffil{\ruipm.

\and

\rumephi.

\and

\rubauman.

\and

\rursuh.}



\enaffil{\enipm.

\and

\enmephi.

\and

\enbauman.

\and

\enrsuh.}



\ruauthorsdetails{3}{% Количество авторов

\fullauthor{\rupid{53328}{Александр Владимирович Манжиров}}

%\CAuthor % Автор-корреспондент

\orcid{0000-0002-7578-6031} % ORCID ID автора 

\regalia{доктор физико-математических наук, профессор} 

\position{заместитель директора} 

\dept{лаб. моделирования в МДТТ\affil{1}} 

\position{профессор} 

\dept{каф. высшей математики\affil{2}} 

\position{профессор} 

\dept{каф. прикладной математики\affil{3}} 

\email{manzh@inbox.ru}

\fullauthor{\rupid{133825}{Михаил Николаевич Михин}}

%\orcid{0000-0000-0000-xxxx} % ORCID ID автора 

\\

\regalia{кандидат физико-математических наук} 

\position{доцент} 

\dept{каф. высшей математики\affil{4}} 

\email[br]{mmikhin@inbox.ru}

\fullauthor{\rupid{53045}{Евгений Валерьевич Мурашкин}}

\CAuthor % Автор-корреспондент

\orcid{0000-0002-3267-4742} % ORCID ID автора 

\regalia{кандидат физико-математических наук} 

\position{старший научный сотрудник} 

\dept{лаб. моделирования в МДТТ\affil{1}} 

\position{доцент} 

\dept{каф. высшей математики\affil{2}} 

\position{доцент} 

\dept{каф. прикладной математики\affil{3}} 

\email[br]{murashkin@ipmnet.ru} \vspace{-3mm}

}



\enauthorsdetails{2}{

\fullauthor{\enpid{53328}{Alexander V. Manzhirov}}

\CAuthor % Сorresponding Author

\orcid{0000-0002-7578-6031} % ORCID ID автора 

\regalia{Dr.  Phys. \& Math. Sci., Professor} 

\position{Deputy Director} 

\dept{Lab. of Modelling in Mechanics of Solids\affil{1}}

\position{Professor} 

\dept{Dept. of High Mathematics\affil{2}} 

\position{Professor} 

\dept{Dept. of  Applied Mathematics\affil{3}} 

\email[br]{manzh@inbox.ru}

\fullauthor{\enpid{133825}{Mikhail N. Mikhin}}

%\orcid{0000-0000-0000-xxxx} % ORCID ID avtora 

\\

\regalia{Cand. Phys. \& Math. Sci., Associate Professor} 

\position{Associate Professor} 

\dept{Dept. of High Mathema\-tics\affil{4}}

\email{mmikhin@inbox.ru}

\fullauthor{\enpid{53045}{Evgenii V. Murashkin}}

\orcid{0000-0002-3267-4742} % ORCID ID

\regalia{Cand. Phys. \& Math. Sci.}

\position{Senior Researcher}

\dept{Lab. of Modelling in Mechanics of Solids\affil{1}}

\position{Asso\-ciate Professor} 

\dept{Dept. of High Mathematics\affil{2}} 

\position{Asso\-ciate Professor} 

\dept{Dept. of  Applied Mathe\-matics\affil{3}} 

\email{murashkin@ipmnet.ru}

}



\rukeywords{рост, вал, форма,  кручение, напряженно"=деформированное состояние}

\enkeywords{growth, shaft, shape,  torsion, stress-strain state}



\ruabstract{Рассматривается задача о  кручении растущего вязкоупругого вала жесткими дисками.  

Вал имеет форму кругового цилиндра. 

С~его торцами соединены жесткие диски. 

Изучается процесс непрерывного роста поверхности вала под действием крутящих моментов, приложенных к~дискам. 

Получены и~исследованы парные сумматорные уравнения, отражающие математическое содержание задачи на различных стадиях процесса наращивания. 

Обсуждаются результаты численного анализа и~особенности качественного поведения основных механических характеристик.}



\enabstract{The torsion of a shaft by rigid disks is considered.  The shaft has the form of circular cylinder. Two rigid disks are attached to its end faces. The process of continuous  growth of such shaft under the influence of twisting torques applied to the disks is studied. Dual series equations which reflect the mathematical content of the problem at the different stages of the growing process are derived and solved. Results of the numerical analysis and singularities of the qualitative mechanical behaviour of the fundamental characteristics are discussed.}



\newdate{manzha}{23}{11}{2017}

\date{\displaydate{manzha}}

\newdate{manzhb}{15}{12}{2017}

\revisiondate{\displaydate{manzhb}}

\newdate{manzhc}{18}{12}{2017}

\accepteddate{\displaydate{manzhc}}



\newdate{manzhe}{20}{12}{2017}

\renewcommand{\JournalDataOnline}{\displaydate{manzhe}}



%\ForPeerReview

\makerutitle



\newpage



\Section[n]{Введение}

Современные методы изготовления деталей сложной формы подразумевают разнообразные технологические процессы обработки, в~том числе связанные с удалением части материала "--- резка, фрезеровка, сверление и т.~д. 

Другой класс технологических процессов основан на синтезе деталей путем последовательного добавления материала на подложку или поверхность произвольной формы. 



Использование новых технологий  позволяет создать пространственную деталь весьма сложной формы из широкого класса материала. 

Однако в~процессе роста в теле неминуемо возникают остаточные деформации. 

С~этими деформациями связаны остаточные напряжения и их уровень может вызывать такие явления, как потеря устойчивости и разрушение, причем они могут произойти уже в процессе такого изготовления. 

Таким образом, разработка и~развитие математических моделей и теоретических методов расчета внутренних технологических напряжений и деформаций в изделиях сложной формы является актуальной научной проблемой как с позиций фундаментальных исследований, так и с точки зрения многочисленных приложений.



Фундаментальными вопросами поверхностного роста занимается механика растущих тел (см., например, \cite{Manzh_bib_01,Manzh_bib_02,Manzh_bib_03,Manzh_bib_022,Manzh_bib_033}). 

В~работе \cite{Manzh_bib_01} развита теория поверхностного роста для исследования двумерных задач наращивания деформируемых тел в случае, когда скоростью деформирования поверхности тел за счет нагрузок и натяга приращиваемых элементов можно пренебречь по сравнению со скоростью притока нового материала к этой поверхности. 

Предложены методы решения новых краевых задач, основанные на приведении неклассических задач наращивания вязкоупругих стареющих тел к задачам теории упругости с некоторым параметром, использовании теории аналитических функций для решения последних и восстановлении истинных параметров напряженно"=деформированного состояния тел (см., например, \cite{Manzh_bib_04,Manzh_bib_05,Manzh_bib_06,Manzh_bib_07,Manzh_bib_08}).



Решение каждой конкретной задачи наращивания деформируемых тел представляет собой самостоятельную и трудоемкую проблему (см., например, \cite{Manzh_bib_10,Manzh_bib_11,Manzh_bib_12}.) 

Однако уже сейчас по виду полученных математических соотношений можно предсказать такие присущие растущим телам явления, как возникновение остаточных напряжений после снятия нагрузок, появление в наращиваемом теле поверхностей разрыва напряжений, зависимость напряженно"=деформированного состояния вязкоупругих тел от скорости и способа их изготовления.



Установлено, что в двумерных задачах концентрации напряжений возле отверстий и кручения в готовом теле без учета процесса наращивания максимум интенсивности касательных напряжений достигается на границе тела. В процессе роста максимум интенсивности касательных напряжений может достигаться на границе раздела основного тела и наращиваемой части, на границе готового тела и в произвольной точке наращиваемой части тела.

Полученные результаты могут служить основой при решении прикладных задач расчета деталей и элементов конструкций, изготавливаемых при помощи методик современной технологической  практики.



\newpage



\Section{Постановка краевой задачи кручения}

Рассмотрим (риc.~\ref{Manzh_Fig_1}) достаточно длинный вал (круговой цилиндр) длиной $ 2l $ и радиусом $ b_0 $ (отношение $ l $ к~$ b_0 $ довольно велико), изготовленный из вязкоупругого стареющего материала в~начальный момент времени. 

К~обоим торцам вала приложены плоские круглые диски радиуса $ a <b_0 $. 

В~момент времени $ \tau_0 $ на диски начинают действовать крутящие моменты $ M (t) $, поворачивая  их на угол $ \gamma = 2\alpha (t) $  друг относительно друга. 

Поверхность вала свободна от напряжений (риc. \ref{Manzh_Fig_1}).



В момент времени $\tau_1$ к~поверхности вала начинается приток вещества. 

Вновь присоединяемые элементы также не напряжены, 

а~время их изготовления совпадает со временем изготовления первоначального тела.



Закон роста вала полностью определяется функцией $b(t)$, которая определяет изменение радиуса вала со временем. 

При этом $b(\tau_1)=b_0$.



Рост прекращается в момент времени $\tau_2$. 

В~этот момент радиус вала равен~$b_1$ ($b(\tau_2)=b_1$), а~его боковая поверхность остается свободной от напряжений при $t\geq\tau_2$. 

Для решения контактной задачи растущего тела воспользуемся квазистатическим приближением в отсутствие объемных сил.



\begin{figure}[!b]

\vspace{-3mm}

\centering

\begin{tikzpicture}[>=stealth]

%Заливка

%цилиндр внутренний

%Нижний эллипс

%\fill[top color=gray!50!black,bottom color=gray!10,middle color=gray,shading=axis,opacity=0.25] (0,0) circle (2cm and 0.5cm);

%боковая поверхность

\fill[left color=gray!50!black,right color=gray!50!black,middle color=gray!50,shading=axis,opacity=0.25] (2,0) -- (2,6) arc (360:180:2cm and .5cm) -- (-2,0) arc (180:360:2cm and 0.5cm);

%верхний эллипс

\fill[top color=gray!90!,bottom color=gray!2,middle color=gray!30,shading=axis,opacity=0.25] (0,6) circle (2cm and 0.5cm);

%цилиндр внешний

\draw[thick] (0,0) circle (4cm and 1cm);

\draw[thick] (-4,0) -- (-4,6);

\draw[thick] (4,0) -- (4,6);

\draw[thick] (0,6) circle (4cm and 1cm);

\draw[thick, densely dashed] (-4,0) arc (180:0:4cm and 1cm);

%цилиндр роста

%нижний эллипс

%\fill[top color=gray!50!black,bottom color=gray!10,middle color=gray,shading=axis,opacity=0.25] (0,0) circle (3cm and 0.75cm);

%верхний эллипс

\fill[top color=gray!50!,bottom color=gray!2,middle color=gray!10,shading=axis,opacity=0.25] (0,6) circle (3cm and 0.75cm);

%боковая поверхность

\fill[left color=gray!30!black,right color=gray!50!black,middle color=gray!50,shading=axis,opacity=0.25] (3,0) -- (3,6) arc (360:180:3cm and .75cm) -- (-3,0) arc (180:360:3cm and 0.75cm);

%контуры

%Ростовой цилиндр

\draw[thick] (-2,6) -- (-2,0) arc (180:360:2cm and 0.5cm) -- (2,6) ++ (-2,0) circle (2cm and 0.5cm);

\draw[thick, densely dashed] (-2,0) arc (180:0:2cm and 0.5cm);

%Внешний

\draw[thick, densely dashed] (0,0) circle (3cm and .75cm);

\draw[thick, densely dashed] (-3,0) -- (-3,6);

\draw[thick, densely dashed] (3,0) -- (3,6);

\draw[thick, densely dashed] (0,6) circle (3cm and .75cm);

%диск верхний

\draw[thick] (0,6.25) circle (1cm and .25cm);

\draw[thick] (-1,6) -- (-1,6.25);

\draw[thick] (1,6) -- (1,6.25);

\draw[thick, densely dashed] (-1,6) arc (180:0:1cm and 0.25cm);

\draw[thick] (-1,6) arc (180:360:1cm and 0.25cm);

%диск нижний

\draw[thick, densely dashed] (0,0) circle (1cm and .25cm);

\draw[thick, densely dashed] (-1,0) -- (-1,-.25);

\draw[thick, densely dashed] (1,0) -- (1,-.25);

\draw[thick, densely dashed] (0,-.25) circle (1cm and .25cm);

%ось 

\draw[loosely dashdotted, ->] (0,7.5) -- (0,-1.5);

\node at (-0.2,-1.3) {$z$};

%указатели радиусов

\draw[->] (0,2) -- (2,2) node[near start,sloped,above] {$b_0$};

\draw[->] (0,3) -- (3,3) node[near start,sloped,above] {$b(t)$};

\draw[->] (0,4) -- (4,4) node[near start,sloped,above] {$b_1$};

%моменты

\node [fill=gray!67,rounded corners=2pt,inner sep=0pt,draw=gray!67] at (0,6.25) {$\small M(t)$};

\node [fill=gray!80,rounded corners=2pt,inner sep=0pt,draw=gray!80] at (0,0) {\small $ M(t)$};

\draw[thick, ->] (-0.35,6.4) arc (120:420:0.7cm and 0.15cm);

\draw[thick, <-] (-0.35,0.15) arc (120:420:0.7cm and 0.15cm);

%длина

\draw (-4.5,6) -- (-4.1,6);

\draw (-4.5,0) -- (-4.1,0);

\draw[<->] (-4.3,0) -- (-4.3,6);

\node [fill=white,draw=white] at (-4.3,3) {$2l$}; 





\end{tikzpicture}

\footnotesize 



\caption{Растущий вал в условиях кручения \label{Manzh_Fig_1}}



[Figure~\ref{Manzh_Fig_1}. Growing Shaft under Torsion]



\end{figure}



Будем считать вал достаточно длинным как во время процесса роста, так и после его прекращения (в этом случае отношения $l/b(t)$ и $l/b_1$ должны быть довольно большими). Принимая во внимание симметричность постановки краевой задачи, рассмотрим только половину вала с одной торцевой поверхностью, зажатой на жестком основании, а c другой торцевой поверхностью, связанной с диском.



Воспользуемся общим подходом для решения подобных задач механики растущих тел, разработанным и описанным в работах \cite{Manzh_bib_01,Manzh_bib_02,Manzh_bib_03,Manzh_bib_022,Manzh_bib_033,Manzh_bib_04,Manzh_bib_05,Manzh_bib_06,Manzh_bib_07,Manzh_bib_08}.



Рассмотрим систему фундаментальных соотношений краевой задачи в интервале времени $t\in[\tau_0,\tau_1]$, тогда для начального вязкоупругого стареющего вала получим

\begin{align}\label{Manzh_eq_1}

  &%(1)

  \frac{\partial\sigma_{r\varphi}}{\partial r}+\frac{\partial\sigma_{\varphi z}}{\partial z}+\frac{2\sigma_{r\varphi}}r=0\quad   (\nabla\cdot\mathbf T= 0),

\\[3mm]

  &\nonumber

  z=0,\ 0\leq r\leq a \ \ \colon  u_\varphi=\alpha(t)r,

\quad 

  z=0,\ a\leq r\leq b_0 \ \colon  \sigma_{\varphi z}=0,

\\

  &\nonumber

  r=b_0,\ 0\leq z\leq l \ \ \! \colon  \sigma_{r\varphi}=0,

\quad  \quad \ \ \!

  z=l,\ \! \ \ 0\leq r\leq b_0\colon u_\varphi=0;

\\[3mm]

  &\nonumber

  \varepsilon_{r\varphi}=\frac12\Bigl(\frac{\partial u_\varphi}{\partial r}-\frac{u_\varphi}r\Bigr),

\quad   \varepsilon_{\varphi z}=\frac12\frac{\partial u_\varphi}{\partial z}\quad

  \Bigl( \mathbf E =\frac12[\nabla\mathbf u+(\nabla\mathbf u)^\top]\Bigr);

\\[3mm]

  &\nonumber

  \mathbf T=2G(t)(\mathbf I+\mathbf L(\tau_0,t))\mathbf E,

\quad 

  (\mathbf I-\mathbf L(\tau_0,t))=(\mathbf I+\mathbf N(\tau_0,t))^{-1};

\\[3mm]

  &\nonumber

   \mathbf L(\tau_0,t)f(t)=\int_{\tau_0}^tf(\tau)K_1(t,\tau)\,d\tau,

\quad 

  K_1(t,\tau)=G(\tau)\frac{\partial}{\partial\tau}\Bigl[\frac1{G(\tau)}+\omega(t,\tau)\Bigr],

\end{align}

где $\mathbf T$ и $\mathbf E$ "--- тензоры напряжений и деформаций  с ненулевыми компонентами $\sigma_{r\varphi}$, $\sigma_{\varphi z}$ и $\varepsilon_{r\varphi}$, $\varepsilon_{\varphi z}$ соответственно; 

$\mathbf u$ "--- вектор перемещений с единственной ненулевой компонентой $u_\varphi$; 

$K_1(t,\tau)$, $\omega(t,\tau)$, $G(t)$ "--- ядро ползучести, мера ползучести и модуль упругой деформации при чистом сдвиге.



Положим 

\begin{equation}%(2)

\label{Manzh_eq_2}

  \mathbf T^\circ=(\mathbf I-\mathbf L(\tau_0,t))\mathbf T G^{-1}

\end{equation}

и подействуем на выражение \eqref{Manzh_eq_1}, содержащее $\mathbf T$ и его компоненты, оператором $(\mathbf I-\mathbf L(\tau_0,t))$. Тогда, учитывая \eqref{Manzh_eq_2}, получим следующую краевую задачу:

\begin{align}

  &\label{Manzh_eq_3}%(3)

  \frac{\partial\sigma_{r\varphi}^\circ}{\partial r}+\frac{\partial\sigma_{\varphi z}^\circ}{\partial z}+\frac{2\sigma_{r\varphi}^\circ}r=0\quad

  (\nabla\cdot\mathbf T^\circ=\mathbf 0),

\\[2mm]

  &\nonumber

  z=0,\ 0\leq r\leq a  \ \ \colon u_\varphi=\alpha(t)r,

\quad 

  z=0,\ a\leq r\leq b_0 \ \colon\sigma_{\varphi z}^\circ=0,

\\

  &\nonumber

  r=b_0,\ 0\leq z\leq l \ \ \! \colon\sigma_{r\varphi}^\circ=0,

\quad  \quad  \ \ \!

  z=l,\ \ \ \! 0\leq r\leq b_0\colon u_\varphi=0;

\\[2mm]

  &\nonumber

  \mathbf E=\frac12[\nabla\mathbf u+(\nabla\mathbf u)^\top],\quad

  \mathbf T^\circ=2\mathbf E.

\end{align}



Исходя из \eqref{Manzh_eq_3} можно получить, что перемещение $u_\varphi$ удовлетворяет уравнению

\begin{equation}\label{Manzh_eq_4}%(4)

  \boldsymbol{\mathcal V}u_\varphi \equiv \frac{\partial^2u_\varphi}{\partial r^2}+\frac{\partial^2u_\varphi}{\partial z^2}+\frac1r\frac{\partial u_\varphi}{\partial r}-\frac{u_\varphi}{r^2}=0.

\end{equation}



Следуя \cite{Manzh_bib_10}, возьмем решение \eqref{Manzh_eq_4} в форме (см. \cite{Manzh_bib_04})

\begin{align}\label{Manzh_eq_5}%(5)

  & 

  u_\varphi(r,z,t)=\frac{ld_0(t)r}{b_0}\Bigl(1-\frac zl\Bigr)

  +\sum_{n=1}^\infty\frac{d_n(t)}{\delta_n}J_1(r\delta_n)\frac{\sh[\delta_n(l-z)]}{\sh(\delta_nl)},

\end{align}

где $d_k(t)$ ($k=0,1,2,\dots$)  "--- известные функции времени, $\delta_n$ "--- неопределенные константы, \linebreak а~$J_\nu(x)$ "--- функции Бесселя порядка $\nu$.



Заметим, что выражение \eqref{Manzh_eq_5} для перемещения $u_\varphi$ удовлетворяет граничному условию в \eqref{Manzh_eq_3} на зажатой торцевой поверхности вала для $z=l$ и позволяет выписать  компоненты тензора преобразованных напряжений $\mathbf T^\circ$ в виде (см. \eqref{Manzh_eq_1} и \eqref{Manzh_eq_3})

\begin{align}\label{Manzh_eq_6}%(6)

\begin{aligned}

  &

  \sigma_{\varphi z}^\circ(r,z,t) = - \frac{d_0(t)r}{b_0}

  -\sum_{n=1}^\infty\!d_n(t)J_1(r\delta_n)\frac{\ch[\delta_n(l\!-\!z)]}{\sh(\delta_nl)},

\\

  &

  \sigma_{r\varphi}^\circ(r,z,t)

  =-\sum_{n=1}^\infty d_n(t)J_2(r\delta_n)\frac{\sh[\delta_n(l-z)]}{\sh(\delta_nl)}.

\end{aligned}

\end{align}



Используя граничное условие из системы \eqref{Manzh_eq_3} на боковой поверхности вала ($r=b_0$) и условия \eqref{Manzh_eq_6}, найдем неизвестные постоянные $\delta_n$. Действительно, приравнивая выражение для

$\sigma_{r\varphi}^\circ$ к нулю при $r=b_0$, получим, что $\delta_n=\lambda_n/b_0$, где $\lambda_n$ "--- корни уравнения $J_2(\lambda_n)=0$.



Наконец, удовлетворяя граничным условиям при $z=0$, для поиска последовательности функций $d_k(t)$ получим следующие парные сумматорные уравнения: 

\begin{equation}\label{Manzh_eq_7}%(7)

\begin{aligned}

  &

  \frac{ld_0(t)r}{b_0}

  +\sum_{n=1}^\infty\frac{b_0d_n(t)}{\lambda_n}J_1\Bigl(\frac{\lambda_nr}{b_0}\Bigr)=\alpha(t)r,

\quad    0\leq r\leq a,

\\

  &

  \frac{d_0(t)r}{b_0}

  +\sum_{n=1}^\infty d_n(t)J_1\Bigl(\frac{\lambda_nr}{b_0}\Bigr)\cth\frac{\lambda_nl}{b_0}=0.

\quad   a\leq r\leq b,

\end{aligned}

\end{equation}



Поскольку $\lambda_n\geq\lambda_1 \approx 3.8317$ и $l/b_0=\varkappa_0\gg1$, выражение для $\cth(\lambda_nl/b_0)$ может быть задано равным единице с высокой степенью точности, и~уравнения \eqref{Manzh_eq_7} можно исследовать в~виде (см. \cite{Manzh_bib_10})

\begin{align}\label{Manzh_eq_8}%(8)

\hskip-3em

\begin{aligned}

  &

  u_\varphi(r,0,t)=\varkappa_0d_0(t)r

  +\sum_{n=1}^\infty \frac{b_0d_n(t)}{\lambda_n}J_1\Bigl(\frac{\lambda_nr}{b_0}\Bigr)=\alpha(t)r,

\quad    0\leq r\leq a,

\\

  &

  \sigma_{\varphi z}^\circ(r,0,t)=\frac{d_0(t)r}{b_0}

  +\sum_{n=1}^\infty d_n(t)J_1\Bigl(\frac{\lambda_nr}{b_0}\Bigr)=0,

\quad  \qquad  \quad  \

  a\leq r\leq b_0.

\end{aligned}

\hskip-3em

\end{align}



Ряды в уравнениях \eqref{Manzh_eq_8} описывают сформулированную контактную задачу в интервале $t\in[\tau_0,\tau_1]$, в который время входит параметрически. Построим теперь решение \eqref{Manzh_eq_8}, получив сначала разрешающие уравнения задачи при непрерывном росте и после прекращения роста. Отметим лишь, что истинные напряжения можно восстановить по формуле

\begin{equation}\label{Manzh_eq_9}%(9)

  \mathbf T(r,z,t) = G(t)\biggl[\mathbf T^\circ(r,z,t)+ \int_{\tau_0}^t  \mathbf T^\circ(r,z,\tau)R_1(t,\tau)\,d\tau\biggr],

\end{equation}

где $R_1(t,\tau)$ "--- резольвента ядра $K_1(t,\tau)$.



Пусть $t\in[\tau_1,\tau_2]$. Тогда краевая задача для растущего вала, скрученного диском, примет вид (см. \cite{Manzh_bib_05,Manzh_bib_06,Manzh_bib_07,Manzh_bib_08})

\begin{align}\label{Manzh_eq_10}%(10)

  &

  \frac{\partial s_{r\varphi}}{\partial r}+\frac{\partial s_{\varphi z}}{\partial z}+\frac{2s_{r\varphi}}r=0\quad

  (\nabla\cdot\mathbf S=\mathbf 0),

\\[2mm]

  &\nonumber

  z=0,\  \ 0\leq r\leq a  \ \ \colon v_\varphi=\dot\alpha(t)r,

\quad \quad \ \quad \quad 

  z=0,\ a\leq r\leq b(t)\colon s_{\varphi z}=0,

\\

  &\nonumber

  r=b(t),\ 0\leq z\leq l \! \ \colon s_{r\varphi}=0,\   t=\tau^*(r),

\quad 

  z=l,\ 0\leq r\leq b(t) \ \! \colon v_\varphi=0;

\\[2mm]

  &\nonumber

  \mathbf D=\frac12[\nabla\mathbf v+(\nabla\mathbf v)^\top],

\quad 

  \mathbf S=2\mathbf D,\quad   \frac{\partial\mathbf T^\circ}{\partial t}=\mathbf S.

\end{align}



Видно, что скорость перемещения $v_\varphi$ удовлетворяет уравнению ${\boldsymbol{\mathcal V}v_\varphi=0}$ (см. \eqref{Manzh_eq_4}), в то время как выражение для $v_\varphi$ и скорость преобразованных нап\-ряжений $s_{r\varphi}$ и $s_{\varphi z}$ могут быть записаны в виде

\begin{equation}\label{Manzh_eq_11}%(11)

\begin{aligned}

  &

  v_\varphi(r,z,t)=\frac{ld_0^\circ(t)r}{b(t)}\Bigl(1-\frac zl\Bigr)

  +\sum_{n=1}^\infty\frac{d_n^\circ(t)}{\eta_n(t)}J_1[r\eta_n(t)]

  \frac{\sh[\eta_n(t)(l-z)]}{\sh[\eta_n(t)l]},

\\

  &

  s_{\varphi z}(r,z,t)=-\frac{d_0^\circ(t)r}{b(t)}

  -\sum_{n=1}^\infty d_n^\circ J_1[r\eta_n(t)]\frac{\ch[\eta_n(t)(l-z)]}{\sh[\eta_n(t)l]},

\\

  &

  s_{r\varphi}(r,z,t)=

  -\sum_{n=1}^\infty d_n^\circ J_2[r\eta_n(t)]\frac{\sh[\eta_n(t)(l-z)]}{\sh[\eta_n(t)l]}.

\end{aligned}

\end{equation}

Здесь $d_k^\circ(t)$ и $\eta_n(t)$ "--- функции, подлежащие определению.



Удовлетворяя граничным условиям \eqref{Manzh_eq_10} и учитывая, что $l/b(t)\gg1$, мы приходим к уравнениям с~бесконечными суммами для нахождения $d_k^\circ(t)$:

\begin{align}\label{Manzh_eq_12}%(12)

\hskip-5em

\begin{aligned}

  &

  v_\varphi(r,0,t)=\varkappa(t)d_0^\circ(t)r

  +\sum_{n=1}^\infty \frac{b(t)d_n^\circ(t)}{\lambda_n}J_1\Bigl[\frac{\lambda_nr}{b(t)}\Bigr]=\dot\alpha(t)r,

\quad 

  0\leq r\leq a,

\\

  &

  s_{\varphi z}(r,0,t)=\frac{d_0^\circ(t)r}{b(t)}

  +\sum_{n=1}^\infty d_n^\circ J_1\Bigl[\frac{r\lambda_n}{b(t)}\Bigr]=0,

\quad    a\leq r\leq b(t),

\\

  &

  \eta_n(t)=\frac{\lambda_n}{b(t)},\quad

  \varkappa(t)=\frac{l}{b(t)},\quad

  \tau_1\leq t\leq\tau_2.

\end{aligned}

\hskip-5em

\end{align}



Если $d_k^\circ(t)$ найдены, что означает, что и $\mathbf S$ и $\mathbf v$ также известны, то тензор напряжений $\mathbf T$ и вектор перемещений $\mathbf u$ можно восстановить по формулам

\begin{align}\label{Manzh_eq_13}%(13)

\hskip-5em

\begin{aligned}

  &

  \mathbf T(r,z,t)=G(t)\biggl\{\frac{\mathbf T(r,z,\tau_0(r))}{G(\tau_0(r))}\biggl[1+\int_{\tau_0(r)}^t  R_1(t,\tau)\,d\tau\biggr] +

\\

  &

  \hskip3em

  + \int_{\tau_0(r)}^t \biggl[\mathbf S(r,z,\tau) + \int_{\tau_0(r)}^\tau  \mathbf S(r,z,\zeta)\,d\zeta R_1(t,\tau)\biggr]\,d\tau\biggr\},

\\

  &

  \mathbf u(r,z,t)=\mathbf u(r,z,\tau_0(r))+\int_{\tau_0(r)}^t\mathbf v(r,z,\tau)\,d\tau.

\end{aligned}

\hskip-5em

\end{align}



Краевая задача для растущего вала, когда прекращается его рост \linebreak $t\geq\tau_2=\tau^*(b_1)$ ($b(t)=b_1$) и выполняются краевые условия $\tau_{r\varphi}=0$, заданные на поверхности вала, имеет вид \eqref{Manzh_eq_10}. 

Как и раньше, ее можно свести к краевой задаче в скоростях перемещения и преобразованных напряжений с решением в виде \eqref{Manzh_eq_11} при условии $b(t)=b_1$. Результирующие сумматорные уравнения сохраняют вид \eqref{Manzh_eq_12}, где $b(t)=b_1$, $\varkappa(t)=\varkappa_1=l/b_1$,

$\eta_n(t)=\eta_n=\lambda_n/b_1$, $t\geq\tau_2$. 

После их решения напряжения $\sigma_{r\varphi}$, $\sigma_{\varphi z}$ и перемещения $u_\varphi$ определяются согласно \eqref{Manzh_eq_13}. 

Следует отметить, что зависимость $\mathbf S$ и $\mathbf v$ от времени $t$ параметрическая.



К полученным уравнениям %в рядах 

должно быть быть добавлено условие равновесия диска

\begin{equation}\label{Manzh_eq_14}%(14)

  M(t)=-2\pi\int_0^a\sigma(\rho,t)\rho^2\,d\rho,\quad

  \sigma(\rho,t)=\sigma_{\varphi z}(\rho,0,t),

\end{equation}

которое выполняется во всем рассматриваемом временном промежутке.



На основании условий \eqref{Manzh_eq_14} могут быть также получены следующие соотношения:

\begin{align}

  &\label{Manzh_eq_18}%(18)

    M^\circ(t)=(\mathbf I-\mathbf L(\tau_0,t))\frac{M(t)}{G(t)}

    =-2\pi\int_0^a\sigma^\circ(\rho,t)\rho^2\,d\rho,

  \quad      \tau_0\leq t\leq\tau_1,

\\

    &

    \frac{\partial M^{\circ}(t)}{\partial t}=\frac1{G(t)}\frac{\partial M(t)}{\partial t}+\int_{\tau_0}^t\frac{\partial M(\tau)}{\partial\tau}\frac{\partial\omega(t,\tau)}{\partial t}d\tau +

  \\

  &\label{Manzh_eq_19}%(19)

    \hskip7em

    +M(\tau_0)\frac{\partial\omega(t,\tau_0)}{\partial t}

    =-2\pi\int_0^as(\rho,t)\rho^2\,d\rho,

  \quad      t\geq\tau_1,

\end{align}

которые в ряде случаев более удобны для построения решения контактной задачи.



\smallskip



\Section{Решение задачи кручения}

Разрешающие парные сумматорные уравнения краевой задачи в трех основных временных интервалах  могут быть представлены общими соотношениями

\begin{equation}\label{Manzh_eq_20}%(20)

\begin{aligned}

  &

  \zeta\varphi_0x+\sum_{n=1}^\infty\frac{\varphi_n}{\lambda_n}J_1(\lambda_nx)=\psi x, \qquad \qquad 

  0\leq x\leq c,

\\

  &

  -p(x)=\varphi_0x+\sum_{n=1}^\infty\varphi_nJ_1(\lambda_nx)=0,\quad

  c\leq x\leq1.

\end{aligned}

\end{equation}

Здесь для $t\in[\tau_0,\tau_1]$ мы положили $\zeta=\varkappa_0$, $\varphi_k=d_k(t)$, $p(x)=\sigma^\circ(xb_0,t)$, $\psi=\alpha(t)$, $c=a/b_0$, $x=r/b_0$;

для $t\in[\tau_1,\tau_2]$  мы имеем  $\zeta=\varkappa(t)$, $\varphi_k=d_k^\circ(t)$, $\psi=\dot\alpha(t)$, $p(x)=s(xb(t),t)$, $c=a/b(t)$, $x=r/b(t)$,  

и для $t\geq\tau_2$ в отличие от предыдущего случая "--- $\zeta=\varkappa_1$, $b(t)=b_1$.



Построим решение \eqref{Manzh_eq_20}, следуя описанному в \cite{Manzh_bib_11} методу. Пусть

\begin{equation}\label{Manzh_eq_21}%(21)

  p(x)=\bigg[\frac{\partial}{\partial x}\int_x^c\frac{g(\xi)}{\sqrt{\xi^2-x^2}}\,d\xi\bigg]h(c-x),

\end{equation}

где $h(c-x)$ "--- функция Хевисайда.



Ряд во втором уравнении \eqref{Manzh_eq_20} является разложением Дини \cite{Manzh_bib_12} функции $(-p(x))$, коэффициенты $\varphi_k$  которого заданы формулами

\begin{align}\label{Manzh_eq_22}%(22)

\hskip-4em

\begin{aligned}

  &

  \varphi_0=-4\int_0^1x^2p(x)\,dx=8\int_0^c\xi g(\xi)\,d\xi,

\\

  &

  \varphi_n=-\frac{2}{J_1^2(\lambda_n)}\int_0^1xp(x)J_1(\lambda_nx)\,dx =

\\

  &

  \hskip2em

  =\frac{2}{J_1^2(\lambda_n)}\int_0^cg(\xi)\sin(\lambda_n\xi)\,d\xi,\quad

  n=1, 2,\ldots,

\end{aligned}

\hskip-4em

\end{align}

в случае, когда учтено условие \eqref{Manzh_eq_21}.



Подставляя соотношения \eqref{Manzh_eq_22} в первое уравнение \eqref{Manzh_eq_20} и используя технику описанную в работах \cite{Manzh_bib_11,Manzh_bib_13,Manzh_bib_14,Manzh_bib_15}, получим интегральное уравнение Фредгольма второго рода для определения функции $g(x)$:

\begin{align}\label{Manzh_eq_23}%(23)

&

  g(x)+\int_0^cg(\xi)k(x,\xi)\,d\xi=\frac{3\psi x}\pi,\qquad

  1\leq x\leq c,

\\

  &

  k(x,\xi)=\frac{16}\pi(1-2\zeta)x\xi

  +\frac4{\pi^2}\int_0^\infty\frac{K_2(y)}{I_2(y)}\bigl[8x\xi I_2(y)-\sh(xy)\sh(\xi y)\bigr]\,dy,

\end{align}

где $K_\nu(y)$, $I_\nu(y)$ "--- функции Бесселя мнимого аргумента порядка $\nu$.



Решение \eqref{Manzh_eq_20}, очевидно, также дает полное решение рассматриваемой контактной задачи. 

Его можно найти с помощью методов, рассмотренных в работах \cite{Manzh_bib_16,Manzh_bib_17}. Здесь же для построения приближенного решения \eqref{Manzh_eq_23} воспользуемся методом, предложенным в \cite{Manzh_bib_10}. Заметим, что для $\zeta\geq10$ отклонение приближенного решения от численного решения не превышает 8.5\,\% при $c=0.7$,  7\,\% при $c=0.6$ и 1\,\% при $c\leq0.5$.



Воспользуемся тем фактом, что величина $\zeta$ достаточно велика, и ограничимся первым членом в выражении для ядра $k(x,\xi)$ (см. \eqref{Manzh_eq_23}):

\begin{equation}\label{Manzh_eq_24}%(24)

  g(x)+\frac{16}\pi(1-2\zeta)x\!\int_0^c g(\xi)\xi\,d\xi=\frac{4\psi x}\pi\qquad

  (1\leq x\leq c).

\end{equation}

Тогда, подставляя $g(x)=Ax$ в уравнения \eqref{Manzh_eq_24} и определяя константу $A$, в~силу \eqref{Manzh_eq_21} будем иметь 

\begin{equation}\label{Manzh_eq_25}%(25)

  p(x)=-\frac{4\psi}{\pi + 16(2\zeta - 1)c^3 /3}\frac{x}{\sqrt{c^2 - x^2}}\qquad    (1\leq x\leq c).

\end{equation}



Зависимости преобразованных контактных напряжений $\sigma^\circ(r,t)$ и их скоростей $s^\circ(r,t)$ от угла поворота диска $\alpha(t)$ могут быть вычислены по формулам

\begin{align}

  &\label{Manzh_eq_26}%(26)

  \sigma^\circ(r,t)=\alpha(t)W(r,b_0),\quad \ \! \

  \tau_0\leq t\leq\tau_1,

\\

  &\label{Manzh_eq_27}%(27)

  s^\circ(r,t)=\dot\alpha(t)W(r,b(t)),\quad

  \tau_1\leq t\leq\tau_2,

\\

  &\label{Manzh_eq_28}%(28)

  s^\circ(r,t)=\dot\alpha(t)W(r,b_1),\quad \ \ \   t\geq\tau_2,

\\

  &\nonumber

  W(r,\xi)=-\frac4{\pi+16(2l/\xi-1)a^3\!/(3\xi^3)}  \frac r{\sqrt{c^2-r^2}}.

\end{align}



Для заданного угла поворота диска контактные напряжения $\sigma^\circ(r,t)$ и их скорости $s^\circ(r,t)$ тут же могут быть найдены из соотношений \eqref{Manzh_eq_26}--\eqref{Manzh_eq_28}, а контактные напряжения $\sigma(r,t)$ могут быть восстановлены по описанным ранее зависимостям. Момент, действующий на диск, вычисляется согласно \eqref{Manzh_eq_14}. Заметим, что при $\alpha(t)=\mathrm{const}$ взаимное влияние начального вала и его вновь образующейся ненапряженной части не проявляется. На основании соотношений \eqref{Manzh_eq_9}, \eqref{Manzh_eq_18} и \eqref{Manzh_eq_26} для заданного крутящего момента $M(t)$ получим

\begin{align}

  &

  \label{Manzh_eq_29}%(29)

  \sigma(r,t)=\frac{3M(t)}{4\pi a^3}\frac r{\sqrt{a^2-r^2}},

\ 

  \alpha(t)=B(b_0)(\mathbf I-\mathbf L(\tau_0,t))\frac{M(t)}{G(t)},\ 

  \tau_0\leq t\leq\tau_1, 

\\

  &

  B(\xi)=\frac3{16a^3}+\frac{2l-\xi}{\pi\xi^4}.

\end{align}



\newpage



Воспользовавшись зависимостями \eqref{Manzh_eq_19}, \eqref{Manzh_eq_22}, \eqref{Manzh_eq_27} и \eqref{Manzh_eq_28}, мы наконец получаем соотношение \eqref{Manzh_eq_29} для определения контактных напряжений в промежутке времени $t\geq\tau_1$ и следующие выражения для угла поворота диска: 

\begin{align*}

  &

  \dot\alpha(t)=\frac{\partial M^\circ(t)}{\partial t}B(b(t)),\quad

  \alpha(t)=\alpha(\tau_1)+\int_{\tau_1}^t\dot\alpha(\tau)\,d\tau,

\quad \tau_1\leq t\leq\tau_2,

\\

  &

  \dot\alpha(t)=\frac{\partial M^\circ(t)}{\partial t}B(b_1),\quad

  \alpha(t)=\alpha(\tau_2)+\int_{\tau_1}^t\dot\alpha(\tau)\,d\tau, 

\quad t\geq\tau_2.

\end{align*}



Оказывается, что рост вала при кручении диска моментом сил оказывает незначительное влияние на распределение контактных напряжений, если радиусы диска и вала не очень велики (конкретные отношения приведены выше). 

Однако в~этом же случае появляется существенная зависимость угла поворота диска от момента и от скорости роста, когда вал начинает расти.



Для расчета параметров напряженно-деформированного состояния вала на некотором расстоянии от его краевых поверхностей (в условиях принципа Сен"--~Венана) можно воспользоваться результатами полученными в~работе~\cite{Manzh_bib_18}.



\smallskip



\Section{Численные расчеты}

Рассмотрим контактную краевую задачу для вала, изготовленного из бетона с модулем упруго-мгновенной деформации сдвига

$G(t)=G=\mathrm{const}$ и мерой ползучести при сдвиге в виде~\cite{Manzh_bib_01}

\begin{equation*}

  \omega(t,\tau)=(D_0+Fe^{-\beta\tau})(1-e^{-\gamma(t-\tau)}).

\end{equation*}



Произведем замену переменных согласно формулам

\begin{gather}

  r^*=\frac ra,\quad

  \rho^*=\frac\rho a,\quad

  t^*=\frac t{\tau_0},\quad

  \sigma^*(r^*,t^*)=\frac{\sigma(r,t)}G,

\quad 

  \tau_1^*=\frac{\tau_1}{\tau_0},\quad

  \tau_2^*=\frac{\tau_2}{\tau_0},

  \\  

  M^*(t^*)=\frac{M(t)}{Ga^3},\quad

  \alpha^*(t^*)=\alpha(t), \quad 

  \beta^*=\beta\tau_0,\quad

  \gamma^*=\gamma\tau_0,\\

  b_0^*=\frac{b_0}a,\quad

  b_1^*=\frac{b_1}a,

\quad 

  b^*(t^*)=\frac{b(t)}a,\quad

  l^*=\frac la,\quad

  D_0^*=D_0G,\quad

  F^*=FG

\end{gather}

и, опуская звездочку в обозначениях, зададим следующие значения функций и параметров:

\begin{gather*}

  b_0=\frac1{0.7},\quad

  l=\frac{290}{0.7},\quad

  b(t)=\frac{b_0(t+\tau_2-2\tau_1)}{\tau_2-\tau_1},

\quad 

  b_1=2b_0,\quad

  M(t)=1,\\

  D_0=0.251,\quad

  F=1.818,

\quad 

  \beta=0.31,\quad

  \gamma=0.6,\quad

  \tau_0=10\,\text{дней}.

\end{gather*}



Видно, что во время роста вала его радиус удваивается. Скорость роста постоянна и определяется только временем начала и прекращения роста. Крутящий момент, действующий на диск, не изменяется со временем. Кроме того, отношение длины вала к его радиусу не менее $20 $ в течение всего процесса роста, в то время как отношение радиуса диска и вала не превосходит $0.7 $, т.~е. можно воспользоваться формулами для поиска приближенного решения.





Что касается распределения контактных напряжений, то достаточно вспомнить соотношения \eqref{Manzh_eq_29}, чтобы отметить, что они (распределения) практически не зависят от свойств материала и процесса роста в рассматриваемом случае.



Поведение угла вращения диска в зависимости от контактных напряжений требует более детального анализа.



Кривые на рис. \ref{Manzh_Fig_2} изображают изменение угла поворота $\gamma=2\alpha$ в момент времени $t$ во время первого процесса наращивания. 

В~этом случае рост вала начинается одновременно с наличием крутящего момента ($\tau_1=1$) для различных скоростей роста $\dot b(t)$: $\dot b(t)=b_0/9$ ($\tau_2=10$) "--- сплошная линия $\dot b(t)=b_0/3$ ($\tau_2=4$) --- штрих-пунктирная линия, и $\dot b(t)=b_0$ ($\tau_2=2$) "--- пунктирная линия. Моменты остановки процесса роста отмечены вертикальными сплошными линиями.





Кривые на рис. \ref{Manzh_Fig_3} соответствуют зависимостям угла вращения диска от времени для случая повторного процесса наращивания. В этом случае для вала, загруженного в момент времени $t=1$, и момента начала роста "--- $\tau_1=2$, нами выбраны различные скорости роста $\dot b(t)$: $\dot b(t)=b_0/8$ ($\tau_2=10$) "--- штрих-пунктирная линия, $\dot b(t)=b_0/2$ ($\tau_2=4$) "--- пунктирная линия. 

Для сравнения, изменение угла поворота диска, скручивающего вал фиксированного радиуса $b_0$, показано сплошной линией. Дуги кривых, расположенные между вертикальными сплошными линиями, характеризуют поведение угла вращения диска в промежутках непрерывного роста валов.









\Section[N]{Выводы}



\begin{enumerate}



\item[1.] Контактные напряжения, действующие на вал из-за взаимодействия с~жесткими дисками, практически не зависят от свойств материала и процесса роста, в случае, когда вал достаточно длинный (отношение длины вала к его радиусу больше или равно 20) и в то же время диски не слишком велики (отношение радиуса диска и вала не превышает $0.7 $) в~течение всего процесса наращивания.

\item[2.] Угол поворота $\gamma(t)$ существенно зависит от скорости роста. Таким образом, предельное значение приращения угла поворота диска $\Delta(\infty)$ ($\Delta(t)=\gamma(t)-\gamma(\tau_0)$) во время медленного роста валов может превышать такое же значение для быстрого роста в $5$ раз и более.

\item[3.]Предельное значение угла поворота диска существенно зависит от временного интервала между началом загрузки и началом роста.

\item[4.]Для постоянного крутящего момента существует характерное время, начиная с которого влиянием процесса наращивания на напряженно-деформированное состояние вязкоупругого вала можно пренебречь.



\end{enumerate}











\AdditionalContent{Конкурирующие интересы}{Заявляем, что в отношении авторства и публикации

этой статьи конфликта интересов не имеем.} 

\AdditionalContent{Авторский вклад и ответственность}{Все  авторы  принимали  участие  в  разработке  концепции  \mbox{статьи}  и  в  написании  рукописи. Авторы несут  полную  ответственность  за  предоставление  окончательной рукописи в печать. Окончательная  версия рукописи была одобрена всеми авторами.} 

\AdditionalContent{Финансирование}{Работа выполнена при частичной финансовой поддержке Российского фонда фундаментальных исследований (проект №~17--01--00712\_a).}



\newpage





\begin{figure}[t!]



\vspace{2mm}



\centering

        \begin{tikzpicture}[>=stealth]

            \begin{axis}[

           xmin=1,

           xmax=13,

           ymin=6,

           ymax=18,

                smooth, height=8cm, width=13cm, grid=major,

               xtick={1,2,4,6,8,10,12,13},

               ytick={6,10,14,18},

               xticklabels={1,2,4,6,8,10,12,$t$},

               yticklabels={6,10,14,$\gamma$},

                xlabel={},

                ylabel={},

                 legend style={at={(0.97,0.3)},

                anchor=north east} ]

           \addplot[color=red, very thick] table[x=xx,y=yy] {./00PIC/Fig_2_1.txt};

           \addplot[color=blue, loosely dashdotted, very thick] table[x=xx,y=yy] {./00PIC/Fig_2_2.txt};

           \addplot[color=green!70!black,dashed, very thick] table[x=xx,y=yy] {./00PIC/Fig_2_3.txt};

                % \legend{$\sigma_{rr}$,$\sigma_{\varphi \varphi}$,$\sigma_{zz}$}

            \draw[->] [thick] (10,0) -- (10,33.5);

            \draw[->] [thick] (30,0) -- (30,79.);

            \draw[->] [thick] (90,0) -- (90,112.);

            \end{axis}

        \end{tikzpicture}

\caption{Угол поворота диска в зависимости от различных скоростей роста в случае первого процесса наращивания 

\label{Manzh_Fig_2}}



\smallskip



{\footnotesize [Figure~\ref{Manzh_Fig_2}. The rotation angle of the disk depending on the different growth rates in the case of the first growth process]}



%\end{figure}

%

\bigskip%

%

%\begin{figure}[!t]

\centering

        \begin{tikzpicture}[>=stealth]

            \begin{axis}[

           xmin=1,

           xmax=13,

           ymin=8,

           ymax=24,

                smooth, height=8cm, width=13cm, grid=major,

               xtick={1,2,4,6,8,10,12,13},

               ytick={8,12,16,20,24},

               xticklabels={1,2,4,6,8,10,12,$t$},

               yticklabels={8,12,16,20,$\gamma$},

                xlabel={},

                ylabel={},

                 legend style={at={(0.97,0.3)},

                anchor=north east} ]

           \addplot[color=red, very thick] table[x=xx,y=yy] {./00PIC/Fig_3_1.txt};

           \addplot[color=blue, loosely dashdotted, very thick] table[x=xx,y=yy] {./00PIC/Fig_3_2.txt};

           \addplot[color=green!70!black,dashed, very thick] table[x=xx,y=yy] {./00PIC/Fig_3_3.txt};

                % \legend{$\sigma_{rr}$,$\sigma_{\varphi \varphi}$,$\sigma_{zz}$}

            \draw[->] [thick] (10,0) -- (10,67.5);

            \draw[->] [thick] (30,0) -- (30,89.5);

            \draw[->] [thick] (90,0) -- (90,149.3);

            \end{axis}

        \end{tikzpicture}

\caption{Угол поворота диска в зависимости от различных скоростей роста в случае процесса повторного наращивания 

\label{Manzh_Fig_3}}



\footnotesize \smallskip



[Figure~\ref{Manzh_Fig_3}. The rotation angle of the disk depending on the different growth rates in the case of the second growth process]



\vspace{-5mm}



\end{figure}







\begin{thebibliography}{99}

  

\RBibitem{Manzh_bib_01}

\by Aрутюнян~Н.~Х., Манжиров~А.~В., Наумов~В.~Э.

\book Контактные задачи механики растущих тел

\yr 1991

\publ Наука

\publaddr М.

\totalpages 176



\RBibitem{Manzh_bib_02}

\by Aрутюнян~Н.~Х., Манжиров~А.~В.

\book Контактные задачи теории ползучести

\yr 1999

\publ Институт механики НАН

\publaddr Ереван

\totalpages 320



\RBibitem{Manzh_bib_03}

\by Манжиров~А.~В.

\paper Общая безынерционная начально-краевая задача для кусочно-непрерывно наращиваемого вязкоупругого стареющего тела

\jour ПММ

\yr 1995

\vol 59

\issue 5

\pages 836--848

%\misctransl

%\by Manzhirov~A.~V.

%\paper The general non-inertial initial-boundary value problem for a viscoelastic ageing solid with piecewise-continuous accretion

%\jour Journal of Applied Mathematics and Mechanics

%\yr 1995

%\vol 59

%\issue 5

%\pages 805--816



\Bibitem{Manzh_bib_022}

\by Manzhirov~A.~V.

\paper Mechanics of growing solids and phase transitions

\jour Key Engineering Materials

\yr 2013

\vol 535--536

\pages 89--93

\crossref{10.4028/www.scientific.net/KEM.535-536.89}



\Bibitem{Manzh_bib_033}

\by Manzhirov~A.~V.

\paper Mechanics of growing solids: New track in mechanical engineering

\inbook ASME 2014 International Mechanical Engineering Congress and Exposition

\bookinfo Volume 9: Mechanics of Solids, Structures and Fluids

\procinfo Montreal, Quebec, Canada, November 14–20, 2014

\yr 2014

\papernumber IMECE2014-36712

\pages V009T12A039

\crossref{10.1115/IMECE2014-36712}

\publaddr ASME



\Bibitem{Manzh_bib_04}

\by Manzhirov~A.~V.

\paper Advances in the Theory of Surface Growth with Applications to Additive Manufacturing Technologies

\jour Procedia Engineering

\yr 2017

\vol 173

\pages 11--16

\crossref{10.1016/j.proeng.2016.12.008}



\Bibitem{Manzh_bib_05}

\by Manzhirov~A.~V., Gupta~N.~K.

\paper Fundamentals of Continuous Growth Processes in Technology and Nature

\jour Procedia IUTAM

\yr 2017

\vol 23

\pages 1--12

\crossref{10.1016/j.piutam.2017.06.001}





\bibitem{Manzh_bib_06}

\by Manzhirov~A.~V.

\paper Mechanical design of viscoelastic parts fabricated using additive manufacturing technologies

\inbook Lecture Notes in Engineering and Computer Science

\bookvol 2

\bookinfo Proceedings of The World Congress on Engineering 2015, WCE 2015

\procinfo  London, U.K., July 1--3, 2015

\yr 2015

\pages 710--714

\publaddr IAENG

\elink \url{http://www.iaeng.org/publication/WCE2015/WCE2015_pp710-714.pdf}





\Bibitem{Manzh_bib_07}

\by Manzhirov~A.~V.

\paper A method for mechanical design of am fabricated viscoelastic parts

\inbook Transactions on Engineering Technologies

\yr 2016

\pages 223--235

\publ Springer

\publaddr Singapore

\crossref{10.1007/978-981-10-1088-0_17}





\Bibitem{Manzh_bib_08}

\by Manzhirov~A.~V.

\paper Fundamentals of mechanical design and analysis for am fabricated parts

\jour Procedia Manufacturing

\yr 2017

\vol 7

\pages 59--65

\crossref{10.1016/j.promfg.2016.12.017}



\Bibitem{Manzh_bib_10}

\by Sneddon~I.~N., Srivastav~R.~P., Mathur~S.~C.

\paper The Reissner-Sagoci problem for a long cylinder of finite radius

\jour Q. J. Mechanics Appl. Math.

\yr 1966

\vol 19

\issue 2

\pages 123--129

\crossref{10.1093/qjmam/19.2.123}



\Bibitem{Manzh_bib_11}

\by Srivastav~R.~P.

\paper Dual series relations. II. Dual relations involving Dini Series

\jour Proc. R. Soc. Edinb. Sec.

\yr 1963

\vol 66

\issue 3

\pages 161--172

\crossref{10.1017/S0080454100007810}



\Bibitem{Manzh_bib_12}

\by Watson~G.~N.

\book A Treatise on the Theory of Bessel Functions.  Second Edition

\yr 1966

\publ Cambridge University Press

\publaddr London

\totalpages vii+804

\zmath{https://zbmath.org/?q=an:0174.36202}





\Bibitem{Manzh_bib_13}

\by Sneddon~I.~N., Srivastav~R.~P.

\paper Dual series relations. I. Dual relations involving Fourier-Bessel series

\jour Proc. R. Soc. Edinb. Sec. A

\yr 1963

\vol 66

\issue 3

\pages 150--160

\crossref{10.1017/S0080454100007809}



\Bibitem{Manzh_bib_14}

\by Srivastav~R.~P.

\paper Dual series relations. III. Dual relations involving trigonometric series

\jour Proc. R. Soc. Edinb. Sec. A

\yr 1963

\vol 66

\issue 3

\pages 173--184

\crossref{10.1017/S0080454100007822}



\Bibitem{Manzh_bib_15}

\by Cooke~J.~C.

\paper The solution of triple and quadruple integral equations and Fourier-Bessel Series

\jour Q. J. Mechanics Appl. Math.

\yr 1972

\vol 25

\issue 2

\pages 247--263

\crossref{10.1093/qjmam/25.2.247}





\Bibitem{Manzh_bib_16}

\by Sneddon~I.~N., Tait~B.~J.

\paper The effect of a penny-shaped crack on the distribution of stress in a long circular

cylinder

\jour Int. J. Eng. Sc.

\yr 1963

\vol 1

\issue 3

\pages 391--409

\crossref{10.1016/0020-7225(63)90016-8}



\Bibitem{Manzh_bib_17}

\by Polyanin~A.~D., Manzhirov~A.~V.

\book Handbook of Integral Equations. Second Edition

\publaddr Boca Raton, London

\publ Chapman \& Hall/ CRC Press

\yr 2008

\totalpages  xxxiv+1108

\mathscinet{http://www.ams.org/mathscinet-getitem?mr=2404728}



\Bibitem{Manzh_bib_18}

\by Manzhirov~A.~V.

\paper Mechanical design of am fabricated prismatic rods under torsion

\jour MATEC Web of Conf.

\yr 2017

\vol 95

\papernumber 12002

\crossref{10.1051/matecconf/20179512002}



\end{thebibliography}





\newpage



\makeentitle



\vspace{-5mm}



\AdditionalContent{Competing interests}{We declare that we have no conflicts of interests with the authorship and publication of this article.} 

\AdditionalContent{Authors' contributions and responsibilities}{Each author has participated in the article concept development and in the manuscript writing. The authors are absolutely responsible for submitting the final manuscript in print. Each author has approved the final version of manuscript.} 

\AdditionalContent{Funding}{This work was partial financially supported by the Russian Foundation for Basic Research (project no. 17--01--00712\_a).}









\begin{thebibliography}{99}

  

\Bibitem{Manzh_bib_01-}

\lang In Russian

\by Arutiunian~N.~Kh., Manzhirov~A.~V., Naumov~V.~E.

\book Kontaktnye zadachi mekhaniki rastushchikh tel \rm [Contact problems of the mechanics of growing bodies]

\yr 1991

\publ Nauka

\publaddr Moscow

\totalpages 176





\Bibitem{Manzh_bib_02-1}

\lang In Russian

\by Arutiunian~N.~Kh., Manzhirov~A.~V.

\book Kontaktnye zadachi teorii polzuchesti \rm [Contact problems of the theory of creep]

\yr 1999

\publ Institute of Mechanics, NAS RA

\publaddr Erevan

\totalpages 320









\Bibitem{Manzh_bib_03-1}

\by Manzhirov~A.~V.

\paper The general non-inertial initial-boundary value problem for a viscoelastic ageing solid with piecewise-continuous accretion

\jour J. Appl. Math. Mech.

\yr 1995

\vol 59

\issue 5

\pages 805--816

\crossref{10.1016/0021-8928(95)00095-X}



\Bibitem{Manzh_bib_0220}

\by Manzhirov~A.~V.

\paper Mechanics of growing solids and phase transitions

\jour Key Engineering Materials

\yr 2013

\vol 535--536

\pages 89--93

\crossref{10.4028/www.scientific.net/KEM.535-536.89}



\Bibitem{Manzh_bib_033-}

\by Manzhirov~A.~V.

\paper Mechanics of growing solids: New track in mechanical engineering

\inbook ASME 2014 International Mechanical Engineering Congress and Exposition

\bookinfo Volume 9: Mechanics of Solids, Structures and Fluids

\procinfo Montreal, Quebec, Canada, November 14–20, 2014

\yr 2014

\papernumber IMECE2014-36712

\pages V009T12A039

\crossref{10.1115/IMECE2014-36712}

\publaddr ASME



\Bibitem{Manzh_bib_04-}

\by Manzhirov~A.~V.

\paper Advances in the Theory of Surface Growth with Applications to Additive Manufacturing Technologies

\jour Procedia Engineering

\yr 2017

\vol 173

\pages 11--16

\crossref{10.1016/j.proeng.2016.12.008}



\Bibitem{Manzh_bib_05=}

\by Manzhirov~A.~V., Gupta~N.~K.

\paper Fundamentals of Continuous Growth Processes in Technology and Nature

\jour Procedia IUTAM

\yr 2017

\vol 23

\pages 1--12

\crossref{10.1016/j.piutam.2017.06.001}





\bibitem{Manzh_bib_06-}

\by Manzhirov~A.~V.

\paper Mechanical design of viscoelastic parts fabricated using additive manufacturing technologies

\inbook Lecture Notes in Engineering and Computer Science

\bookvol 2

\bookinfo Proceedings of The World Congress on Engineering 2015, WCE 2015

\procinfo  London, U.K., July 1--3, 2015

\yr 2015

\pages 710--714

\publaddr IAENG

\elink \url{http://www.iaeng.org/publication/WCE2015/WCE2015_pp710-714.pdf}





\Bibitem{Manzh_bib_07-}

\by Manzhirov~A.~V.

\paper A method for mechanical design of am fabricated viscoelastic parts

\inbook Transactions on Engineering Technologies

\yr 2016

\pages 223--235

\publ Springer

\publaddr Singapore

\crossref{10.1007/978-981-10-1088-0_17}





\Bibitem{Manzh_bib_08-}

\by Manzhirov~A.~V.

\paper Fundamentals of mechanical design and analysis for am fabricated parts

\jour Procedia Manufacturing

\yr 2017

\vol 7

\pages 59--65

\crossref{10.1016/j.promfg.2016.12.017}



\Bibitem{Manzh_bib_10-}

\by Sneddon~I.~N., Srivastav~R.~P., Mathur~S.~C.

\paper The Reissner-Sagoci problem for a long cylinder of finite radius

\jour Q. J. Mechanics Appl. Math.

\yr 1966

\vol 19

\issue 2

\pages 123--129

\crossref{10.1093/qjmam/19.2.123}



\Bibitem{Manzh_bib_11-}

\by Srivastav~R.~P.

\paper Dual series relations. II. Dual relations involving Dini Series

\jour Proc. R. Soc. Edinb. Sec.

\yr 1963

\vol 66

\issue 3

\pages 161--172

\crossref{10.1017/S0080454100007810}



\Bibitem{Manzh_bib_12-1}

\by Watson~G.~N.

\book A Treatise on the Theory of Bessel Functions.  Second Edition

\yr 1966

\publ Cambridge University Press

\publaddr London

\totalpages vii+804

\zmath{https://zbmath.org/?q=an:0174.36202}





\Bibitem{Manzh_bib_13-1}

\by Sneddon~I.~N., Srivastav~R.~P.

\paper Dual series relations. I. Dual relations involving Fourier-Bessel series

\jour Proc. R. Soc. Edinb. Sec. A

\yr 1963

\vol 66

\issue 3

\pages 150--160

\crossref{10.1017/S0080454100007809}



\Bibitem{Manzh_bib_14-1}

\by Srivastav~R.~P.

\paper Dual series relations. III. Dual relations involving trigonometric series

\jour Proc. R. Soc. Edinb. Sec. A

\yr 1963

\vol 66

\issue 3

\pages 173--184

\crossref{10.1017/S0080454100007822}



\Bibitem{Manzh_bib_15-1}

\by Cooke~J.~C.

\paper The solution of triple and quadruple integral equations and Fourier-Bessel Series

\jour Q. J. Mechanics Appl. Math.

\yr 1972

\vol 25

\issue 2

\pages 247--263

\crossref{10.1093/qjmam/25.2.247}





\Bibitem{Manzh_bib_16-1}

\by Sneddon~I.~N., Tait~B.~J.

\paper The effect of a penny-shaped crack on the distribution of stress in a long circular

cylinder

\jour Int. J. Eng. Sc.

\yr 1963

\vol 1

\issue 3

\pages 391--409

\crossref{10.1016/0020-7225(63)90016-8}



\Bibitem{Manzh_bib_17-1}

\by Polyanin~A.~D., Manzhirov~A.~V.

\book Handbook of Integral Equations. Second Edition

\publaddr Boca Raton, London

\publ Chapman \& Hall/ CRC Press

\yr 2008

\totalpages  xxxiv+1108

\mathscinet{http://www.ams.org/mathscinet-getitem?mr=2404728}



\Bibitem{Manzh_bib_18-1}

\by Manzhirov~A.~V.

\paper Mechanical design of am fabricated prismatic rods under torsion

\jour MATEC Web of Conf.

\yr 2017

\vol 95

\papernumber 12002

\crossref{10.1051/matecconf/20179512002}



\end{thebibliography}









%\end{document}

